\documentclass[10pt,twocolumn]{article}

% ── Packages ──────────────────────────────────────────────────────────────────
\usepackage[T1]{fontenc}
\usepackage[utf8]{inputenc}
\usepackage{lmodern}
\usepackage[margin=1in,columnsep=0.25in]{geometry}
\usepackage{microtype}
\usepackage{amsmath,amssymb}
\usepackage{booktabs}
\usepackage{multirow}
\usepackage{graphicx}
\usepackage{xcolor}
\usepackage{listings}
\usepackage{enumitem}
\usepackage[hidelinks,
            colorlinks=true,
            linkcolor=blue!60!black,
            citecolor=blue!60!black,
            urlcolor=blue!60!black]{hyperref}
\usepackage{fancyhdr}
\usepackage{titlesec}
\usepackage{abstract}

% ── Typography ────────────────────────────────────────────────────────────────
\titleformat{\section}{\large\bfseries}{\thesection.}{0.5em}{}
\titleformat{\subsection}{\normalsize\bfseries}{\thesubsection.}{0.5em}{}
\setlength{\parskip}{3pt}
\setlength{\parindent}{1em}

% ── Code listings ─────────────────────────────────────────────────────────────
\lstset{
  basicstyle=\ttfamily\footnotesize,
  breaklines=true,
  frame=single,
  framerule=0.4pt,
  columns=flexible,
  keepspaces=true,
  backgroundcolor=\color{gray!8},
}

% ── Header / Footer ───────────────────────────────────────────────────────────
\pagestyle{fancy}
\fancyhf{}
\renewcommand{\headrulewidth}{0.4pt}
\fancyhead[L]{\footnotesize Hybrid ANE-MLX-Bench}
\fancyhead[R]{\footnotesize AtomGradient, 2026}
\fancyfoot[C]{\thepage}

% ── Title ─────────────────────────────────────────────────────────────────────
\title{\textbf{Disaggregated LLM Inference on Apple Silicon:\\
       CoreML ANE Prefill + MLX GPU Decode}}

\author{%
  AtomGradient\\
  \small\href{https://github.com/AtomGradient/hybrid-ane-mlx-bench}{github.com/AtomGradient/hybrid-ane-mlx-bench}
}

\date{2026}

% ══════════════════════════════════════════════════════════════════════════════
\begin{document}
\maketitle

% ── Abstract ──────────────────────────────────────────────────────────────────
\begin{abstract}
We characterize four inference strategies for Qwen3.5 language models on
Apple Silicon, each targeting a different combination of the CPU, GPU, and
Neural Engine (ANE):
(1)~\textbf{MLX baseline} --- GPU-only inference via the MLX framework;
(2)~\textbf{Hybrid CoreML+MLX} --- batched ANE prefill via CoreML with GPU
decode via MLX, bridged through a custom KV-cache converter;
(3)~\textbf{ANE-LM} --- fully sequential ANE inference using private
\texttt{AppleNeuralEngine.framework} APIs;
(4)~\textbf{ANE-LM Hybrid} --- sequential ANE prefill with MLX GPU decode,
bridged via a binary cache transfer (KV, SSM, and conv states).
Experiments on an M2 Ultra (192\,GB, 800\,GB/s) with Qwen3.5-0.8B (FP16)
reveal four key findings:
(i)~ANE batched prefill (CoreML, \texttt{seq512}) reaches parity with GPU
prefill at $\sim$410 prompt tokens;
(ii)~ANE sequential inference (private API) yields a constant dispatch
latency of $\sim$42\,ms/token regardless of prompt length;
(iii)~decode throughput scales linearly with memory bandwidth, confirming
the decode phase is bandwidth-bound rather than compute-bound;
and (iv)~offloading prefill to the ANE reduces GPU power draw from
62.05\,W to 0.22\,W (a 282$\times$ reduction), saving 60.25\,W of
package power --- a critical advantage for mobile devices operating within
strict thermal budgets.
\end{abstract}

% ══════════════════════════════════════════════════════════════════════════════
\section{Introduction}
\label{sec:intro}

Apple Silicon provides three distinct compute units --- CPU, GPU, and Neural
Engine (ANE) --- on a unified memory bus, making it an attractive platform for
edge LLM inference.
Despite this heterogeneity, existing frameworks target only one unit at a
time: MLX routes all operations through the GPU via Metal, while CoreML
dispatches fixed-shape workloads to the ANE.

This work asks: can we improve end-to-end LLM inference latency or
throughput by distributing the prefill and decode phases across different
compute units?
We evaluate two complementary approaches alongside a pure-GPU baseline, and
introduce and fully measure a fourth pipeline combining the best of each.

\textbf{Contribution 1: Hybrid CoreML+MLX pipeline.}
We implement a disaggregated inference pipeline in which the prompt-prefill
phase runs as a single batched forward pass through CoreML (targeting ANE),
while autoregressive decode runs via MLX on the GPU.
The key engineering challenge is bridging the KV cache between frameworks
for Qwen3.5's hybrid attention architecture (DeltaNet + full softmax).

\textbf{Contribution 2: Private ANE API inference (ANE-LM).}
We benchmark ANE-LM~\cite{anelm}, a C++ inference engine that uses the
private \texttt{AppleNeuralEngine.framework} API to dispatch matmul
operations directly to the ANE without the CoreML abstraction layer.
This establishes the single-token ANE dispatch latency floor.

\textbf{Contribution 3: Systematic hardware characterization.}
We benchmark four Qwen3.5 variants (0.8B--9B) across three prompt lengths
on M2 Ultra, providing a comprehensive picture of how model size,
quantization, and prompt length interact with Apple Silicon's compute units.

% ══════════════════════════════════════════════════════════════════════════════
\section{Background}
\label{sec:background}

\subsection{Qwen3.5 Architecture}

Qwen3.5 uses a hybrid per-layer attention design with
\texttt{full\_attention\_interval}$=4$:
\begin{itemize}[topsep=2pt,itemsep=1pt]
  \item \textbf{DeltaNet layers} (3/4 of layers): Linear attention with a
    recurrent state $\mathbf{S} \in \mathbb{R}^{B \times H_v \times D_v \times D_k}$
    and a causal convolution state.
    No KV cache; state is a fixed-size recurrent tensor.
  \item \textbf{Full-attention layers} (1/4 of layers): Scaled dot-product
    attention with output gating, QK-norm, and M-RoPE.
    These layers use a standard KV cache.
\end{itemize}

This heterogeneous cache structure requires a non-trivial bridge layer when
transferring prefill state from CoreML to MLX.

\subsection{Apple Silicon Compute Units}

\begin{table}[h]
\centering
\footnotesize
\begin{tabular}{@{}p{2.1cm}p{2.2cm}p{2.2cm}@{}}
\toprule
Unit & Strengths & Weaknesses \\
\midrule
GPU (MLX/Metal) & Dynamic shapes, lazy eval & Higher per-op overhead \\
ANE (CoreML) & High TOPS, energy eff. & Fixed shapes, fixed seq \\
ANE (private API) & Low framework overhead & Sequential dispatch \\
\bottomrule
\end{tabular}
\caption{Apple Silicon compute units and their inference trade-offs.}
\label{tab:units}
\end{table}

\noindent
On M2 Ultra: 76 GPU cores, 32 ANE cores (31.6 TOPS),
800\,GB/s unified memory bandwidth.

\subsection{Inference Phase Characteristics}

\textbf{Prefill} is compute-bound: processing $N$ prompt tokens in a single
batched forward pass exploits data parallelism across the sequence dimension.
FLOP count scales as $\mathcal{O}(N d^2)$ for attention and
$\mathcal{O}(N d h)$ for FFN layers.

\textbf{Decode} is memory-bandwidth-bound: each autoregressive step reads all
model weights to produce a single token.
Theoretical throughput upper bound:
\begin{equation}
  T_{\max} = \frac{B_{\text{mem}}}{W_{\text{bytes}}}
  \label{eq:bw_bound}
\end{equation}
where $B_{\text{mem}}$ is memory bandwidth (GB/s) and $W_{\text{bytes}}$ is
the model weight size in bytes.

% ══════════════════════════════════════════════════════════════════════════════
\section{System Design}
\label{sec:design}

\subsection{Hybrid CoreML+MLX Pipeline}

\begin{figure}[t]
\begin{lstlisting}
Prompt tokens
     |
     v
CoreML Prefill (ANE)
  Fixed seq_len: 64 / 256 / 512
  Left-pads prompt to seq_len
  Outputs: logits + cache tensors
     |
     v
Cache Bridge
  DeltaNet -> ArraysCache(size=2):
    cache[0] = conv_state
    cache[1] = delta_state
  Full-attn -> KVCache:
    keys, values [B, H, L, D]
    offset = prompt_len (for RoPE)
     |
     v
MLX Decode (GPU)
  Autoregressive, one token/step
\end{lstlisting}
\caption{Hybrid CoreML+MLX pipeline: ANE prefill followed by MLX GPU decode,
  connected via a cache bridge that converts CoreML output tensors to
  mlx\_lm's KVCache/ArraysCache format.}
\label{fig:pipeline}
\end{figure}

\subsection{CoreML Model Conversion}

The ANE-compatible CoreML model is derived from original HuggingFace weights:
\begin{itemize}[topsep=2pt,itemsep=1pt]
  \item \textbf{Linear $\to$ Conv2d}: All linear layers are replaced by
    \texttt{Conv2d(kernel\_size=1)} for ANE compatibility.
  \item \textbf{RMSNorm}: Implemented via the identity
    $\mathbf{x}/\sqrt{\text{mean}(\mathbf{x}^2)+\epsilon}$ using
    \texttt{layer\_norm} with unit scale/zero bias.
  \item \textbf{DeltaNet recurrence}: Unrolled as a functional loop
    (no in-place ops, required by CoreML MIL).
  \item \textbf{Fixed seq\_len}: Traced at compile time for 64, 256, or 512;
    prompts are left-padded to the nearest supported length.
\end{itemize}

On macOS 26.3, \texttt{CPU\_AND\_NE} causes an ANE IPC daemon deadlock;
\texttt{compute\_units=ALL} must be used instead.

\subsection{ANE-LM Private API Engine}

ANE-LM~\cite{anelm} bypasses CoreML and uses the private
\texttt{AppleNeuralEngine.framework} to compile and execute matmul kernels
directly on the ANE:
\begin{enumerate}[topsep=2pt,itemsep=1pt]
  \item Weights are compiled into ANE kernel blobs via
    \texttt{\_ANEInMemoryModelDescriptor} at load time.
  \item Per-token inference: activations are written to
    \texttt{IOSurface} shared memory buffers, then dispatched to ANE.
  \item CPU operations (RMSNorm, RoPE, attention, sampling) execute via
    Accelerate/BLAS.
\end{enumerate}

This one-token-per-dispatch design makes TTFT scale linearly with prompt
length with no batching benefit.

\subsection{ANE-LM Hybrid (Cache Bridge Pipeline)}

The fourth configuration runs \textit{ANE-LM private API prefill} followed
by \textit{MLX GPU decode}, bridging the two frameworks via a binary cache file.

ANE-LM and MLX use different internal cache formats:
ANE-LM stores KV caches as contiguous float32 buffers in
$[\text{N}, H_{\text{kv}}, D]$ layout, SSM recurrent states as
$[H_v, D_k, D_v]$ (transposed relative to MLX's expectation), and conv
states as a circular buffer with a position pointer.
We implement a C++ \texttt{export\_cache()} method in ANE-LM and a Python
bridge that reads the binary file and converts each tensor to the shapes
required by \texttt{mlx\_lm}:
\begin{itemize}[topsep=2pt,itemsep=1pt]
  \item \textbf{KV cache}: $[\text{N}, H_{\text{kv}}, D] \to [1, H_{\text{kv}}, \text{N}, D]$
  \item \textbf{SSM state}: $[H_v, D_k, D_v] \to [1, H_v, D_v, D_k]$
    (swap last two dims)
  \item \textbf{Conv state}: circular buffer $[C, K{-}1] \to [1, K{-}1, C]$
    (unroll circular order, then transpose)
\end{itemize}
The \texttt{ane-lm prefill} subcommand writes the binary file and prints
JSON timing metadata to stdout; the Python engine reads these in sequence.

% ══════════════════════════════════════════════════════════════════════════════
\section{Evaluation}
\label{sec:eval}

\subsection{Experimental Setup}

\begin{itemize}[topsep=2pt,itemsep=1pt]
  \item \textbf{Hardware}: Apple M2 Ultra, 192\,GB unified memory,
    800\,GB/s bandwidth, 32-core ANE (31.6\,TOPS)
  \item \textbf{OS}: macOS 26.3 (25D125)
  \item \textbf{Software}: MLX\,0.31.0, Python\,3.11, coremltools\,8.x;
    ANE-LM built with clang (Release)
  \item \textbf{Methodology}: 4 runs per configuration; first run discarded
    as warmup; median of remaining 3 reported
  \item \textbf{Sampling}: temperature$=0$ (greedy), max 200 generated tokens
\end{itemize}

\textbf{Prompt lengths:}
(1) \textit{short} --- 6 tokens (``Explain neural networks in one sentence.'');
(2) \textit{medium} --- 133 tokens (Apple Silicon UMA paragraph);
(3) \textit{long} --- 410 tokens (edge LLM inference research abstract).
ANE-LM appends a Qwen3.5 chat template, adding $\approx$12 system tokens.

\subsection{Models}

\begin{table}[h]
\centering
\small
\begin{tabular}{@{}llrrc@{}}
\toprule
Model & Quant & Params & Size & CoreML \\
\midrule
Qwen3.5-0.8B & FP16 & 0.85B & 1.6\,GB & \checkmark \\
Qwen3.5-2B & 8-bit & 1.54B & 2.0\,GB & $\times$ \\
Qwen3.5-2B & BF16 & 1.54B & 4.0\,GB & \checkmark \\
Qwen3.5-9B & 8-bit & 9.36B & 9.5\,GB & $\times$ \\
\bottomrule
\end{tabular}
\caption{Models benchmarked. ``CoreML'' indicates FP16/BF16 HF weights
  are available for CoreML conversion.}
\label{tab:models}
\end{table}

\subsection{MLX GPU Baseline}

Table~\ref{tab:mlx_baseline} shows pure-GPU MLX decode throughput across
all four model variants.
Decode throughput scales with the bandwidth-to-weight-bytes ratio
(Eq.~\ref{eq:bw_bound}), confirming bandwidth-bound behavior.

\begin{table}[h]
\centering
\small
\begin{tabular}{@{}llrrrr@{}}
\toprule
Model & Quant & Toks & TTFT & Dec.\ & Mem \\
 & & & (ms) & (tok/s) & (GB) \\
\midrule
\multirow{3}{*}{0.8B} & \multirow{3}{*}{FP16}
  & 6 & 56 & 71.5 & 3.42 \\
  & & 133 & 69 & 70.2 & 3.58 \\
  & & 410 & 96 & 69.0 & 4.18 \\
\midrule
\multirow{3}{*}{2B} & \multirow{3}{*}{8-bit}
  & 6 & 14 & 141.8 & 2.51 \\
  & & 133 & 73 & 141.0 & 2.74 \\
  & & 410 & 162 & 138.9 & 3.23 \\
\midrule
\multirow{3}{*}{2B} & \multirow{3}{*}{BF16}
  & 6 & 22 & 101.3 & 4.16 \\
  & & 133 & 54 & 100.6 & 4.37 \\
  & & 410 & 123 & 99.7 & 4.67 \\
\midrule
\multirow{3}{*}{9B} & \multirow{3}{*}{8-bit}
  & 6 & 39 & 56.4 & 9.76 \\
  & & 133 & 265 & 56.1 & 10.00 \\
  & & 410 & 625 & 56.5 & 10.43 \\
\bottomrule
\end{tabular}
\caption{MLX GPU baseline results. Decode (tok/s) is stable across
  prompt lengths; TTFT grows with tokens.}
\label{tab:mlx_baseline}
\end{table}

\subsection{Four-Pipeline Comparison (Qwen3.5-0.8B)}

Table~\ref{tab:pipeline_ttft} summarises TTFT across all four inference
strategies for Qwen3.5-0.8B FP16.

\begin{table*}[t]
\centering
\small
\begin{tabular}{@{}llrrrr@{}}
\toprule
Pipeline & Prefill & \multicolumn{3}{c}{TTFT (ms)} & Decode \\
\cmidrule(lr){3-5}
 & engine & short (6/18) & med (133/145) & long (410/422) & (tok/s) \\
\midrule
MLX GPU (baseline) & GPU batched & 56 & 69 & 96 & 69--72 \\
Hybrid CoreML+MLX & ANE batched & 274 & 411 & \textbf{100} & 69--73 \\
ANE-LM (private API) & ANE sequential & 769 & 5{,}867 & 17{,}831 & 23--24 \\
ANE-LM Hybrid (bridge) & ANE sequential & 767 & 6{,}060 & 17{,}601 & 67--70 \\
\bottomrule
\end{tabular}
\caption{Time-to-first-token (TTFT) for all four pipelines, Qwen3.5-0.8B FP16, M2 Ultra.
  ANE-LM token counts include $\approx$12 chat-template overhead tokens (shown as
  ``6/18'', ``133/145'', ``410/422'').
  All four rows are \textbf{end-to-end measured} (median of 3 post-warmup runs).
  ANE-LM Hybrid decode speed is measured live via MLX GPU after cache bridge load.}
\label{tab:pipeline_ttft}
\end{table*}

Table~\ref{tab:hybrid_ane} details the Hybrid CoreML+MLX prefill phase.

\begin{table}[h]
\centering
\small
\begin{tabular}{@{}lrrrr@{}}
\toprule
Prompt & Tokens & seq\_len & Prefill & Prefill \\
 & & & (ms) & (tok/s) \\
\midrule
short & 6 & 64 & 274 & 22 \\
medium & 133 & 256 & 410 & 325 \\
long & 410 & 512 & 99 & 4{,}128 \\
\bottomrule
\end{tabular}
\caption{Hybrid CoreML+MLX prefill performance, Qwen3.5-0.8B FP16.
  seq\_len is the smallest available model size (64/256/512) $\geq$ prompt length.
  A 128-token model was not built; medium prompts therefore route to seq256.}
\label{tab:hybrid_ane}
\end{table}

Table~\ref{tab:hybrid_2b} extends the Hybrid CoreML+MLX results to the
larger Qwen3.5-2B BF16 model.

\begin{table}[h]
\centering
\small
\begin{tabular}{@{}lrrrr@{}}
\toprule
Prompt & Tokens & seq\_len & TTFT & Decode \\
 & & & (ms) & (tok/s) \\
\midrule
short & 6 & 64 & 22 & 104.2 \\
medium & 133 & 256 & 54 & 102.3 \\
long & 410 & 512 & 122 & 100.7 \\
\bottomrule
\end{tabular}
\caption{Hybrid CoreML+MLX results for Qwen3.5-2B BF16 (M2 Ultra).
  TTFT matches the MLX GPU baseline at all prompt lengths
  (22\,ms / 54\,ms / 122\,ms vs.\ baseline 22\,ms / 54\,ms / 123\,ms),
  while decode throughput matches baseline (100--104 vs.\ 100--101\,tok/s).}
\label{tab:hybrid_2b}
\end{table}

\textbf{Key finding}: for the 2B BF16 model, Hybrid CoreML+MLX matches
MLX GPU performance at \emph{all} tested prompt lengths, with no dispatch
overhead penalty.
This contrasts with 0.8B, where short/medium prompts were 5--6$\times$
slower.
The likely cause is that 2B's larger hidden dimension (1536 vs.\ 1024)
provides enough compute to amortize CoreML dispatch overhead even at
seq64.

Table~\ref{tab:ane_lm} shows ANE-LM sequential inference.

\begin{table}[h]
\centering
\small
\begin{tabular}{@{}lrrrr@{}}
\toprule
Prompt & Tokens & TTFT & Prefill & Decode \\
 & (w/ tmpl) & (ms) & (tok/s) & (tok/s) \\
\midrule
short & 18 & 769 & 23.4 & 24.3 \\
medium & 145 & 5{,}867 & 24.7 & 23.8 \\
long & 422 & 17{,}831 & 23.7 & 22.8 \\
\bottomrule
\end{tabular}
\caption{ANE-LM private API inference results. Token counts include the
  Qwen3.5 chat template ($\approx$12 tokens overhead).}
\label{tab:ane_lm}
\end{table}

% ══════════════════════════════════════════════════════════════════════════════
\section{Discussion}
\label{sec:discussion}

\subsection{Decode is Memory-Bandwidth Bound}

Decode throughput is proportional to memory bandwidth and inversely
proportional to model weight bytes (Eq.~\ref{eq:bw_bound}):
\begin{itemize}[topsep=2pt,itemsep=1pt]
  \item \textbf{2B 8-bit (139--142\,tok/s)} outperforms \textbf{2B BF16
    (100--101\,tok/s)} by $\approx$40\%: quantization halves bytes read per step.
  \item \textbf{0.8B FP16 (69--72\,tok/s)} is slower than 2B 8-bit despite
    2.5$\times$ fewer parameters: FP16 uses 2\,bytes/element vs.\
    8-bit's 1\,byte/element.
  \item \textbf{9B 8-bit (56--57\,tok/s)} is bottlenecked by raw weight volume
    (9.5\,GB at 800\,GB/s).
\end{itemize}

The theoretical ceiling for 0.8B FP16 from Eq.~\ref{eq:bw_bound} is
$800\,\text{GB/s} / 1.6\,\text{GB} \approx 500\,\text{tok/s}$;
measured 71\,tok/s implies $\approx$14\% bandwidth utilization,
consistent with KV cache reads, normalization, and sampling overhead.

\subsection{Prefill: Batched vs. Sequential ANE Dispatch}

The three prefill strategies differ fundamentally in their compute pattern:

\textbf{MLX GPU (batched)}: processes the full prompt in one forward pass.
GPU prefill throughput scales with prompt length due to increasing GPU
utilization ($107\,\text{tok/s}$ at 6 tokens; $4{,}270\,\text{tok/s}$ at
410 tokens).

\textbf{Hybrid CoreML (batched)}: processes all tokens in one CoreML call
on the ANE.
CoreML dispatch overhead ($\approx$250\,ms for seq64) dominates at short
prompts; at seq512 (410 tokens), ANE achieves $4{,}128\,\text{tok/s}$ ---
97\% of GPU prefill throughput.
\textbf{Crossover}: $\approx$410 prompt tokens (seq512); below this,
CoreML overhead exceeds the GPU prefill time.

\textbf{ANE-LM private API (sequential)}: one token per ANE dispatch.
Single-token dispatch latency is constant at $\approx$42\,ms/token:
\begin{equation}
  \tau_{\text{dispatch}} = \frac{769\,\text{ms}}{18\,\text{tok}}
  \approx 42.7\,\text{ms/token}
  \label{eq:dispatch}
\end{equation}

TTFT scales linearly with prompt length (Figure implied by
Table~\ref{tab:ane_lm}), with zero batching benefit.
ANE-LM sequential achieves 174$\times$ lower prefill throughput than
CoreML batched (23.7 vs.\ 4{,}128\,tok/s).

\subsection{Effect of Replacing ANE Decode with GPU Decode}

The \textit{ANE-LM Hybrid} measured pipeline confirms that replacing ANE
decode (23--24\,tok/s) with GPU decode (67--70\,tok/s) yields a 3$\times$
decode speedup at no measurable additional power cost:
GPU power during decode is 14.17\,W (MLX baseline) vs.\ 12.37\,W
(ANE-LM Hybrid), a difference within measurement variance (see
Section~\ref{sec:power}).
However, for typical conversational prompts ($\leq$422 tokens), prefill
TTFT dominates total latency:
the long-prompt TTFT alone (17{,}601\,ms) is larger than generating
200 tokens at GPU decode speed ($200/70 \times 1000 = 2{,}857\,\text{ms}$)
by a factor of $6\times$.
Any pipeline based on sequential ANE prefill is therefore unsuitable for
interactive inference at these prompt lengths, regardless of decode speed.

\textbf{Implication}: the bottleneck is \emph{batched vs.\ sequential
dispatch}, not ANE hardware throughput.
CoreML batched prefill (Pipeline~2) achieves 4{,}128\,tok/s on the same ANE,
matching GPU throughput at 410 tokens (100\,ms vs.\ 96\,ms).
If the private \texttt{AppleNeuralEngine.framework} API exposed a batched
dispatch interface, ANE-LM Hybrid TTFT would theoretically approach GPU
parity --- but the current API is limited to single-token dispatch.

\subsection{When Does the Hybrid CoreML Approach Help?}

For \textbf{0.8B FP16} on M2 Ultra, Hybrid CoreML+MLX is only beneficial when
prompt length exceeds $\approx$410 tokens (where CoreML overhead becomes
negligible); for shorter prompts, MLX GPU is strictly better.
However, \textbf{2B BF16} shows no dispatch penalty at any tested length
(Table~\ref{tab:hybrid_2b}): TTFT matches GPU baseline even at 6 tokens.
The larger model's compute-to-dispatch ratio amortizes CoreML overhead
at all prompt lengths.
Factors that shift the crossover point:
\begin{itemize}[topsep=2pt,itemsep=1pt]
  \item \textbf{Weaker GPU}: M1 Max (32 cores) has $\approx$2$\times$
    slower GPU prefill, shifting the crossover to $\sim$200--256 tokens.
  \item \textbf{Longer context}: seq\_len $>$ 512 would further amortize
    the CoreML dispatch overhead.
  \item \textbf{GPU contention}: ANE prefill provides hardware isolation
    at no additional TTFT cost when the GPU is occupied.
  \item \textbf{Storage overhead}: Hybrid deployment requires both the
    original HuggingFace weights (4.5\,GB for Qwen3.5-0.8B FP16) and
    the compiled CoreML packages
    (seq64: 1.9\,GB, seq256: 1.9\,GB, seq512: 2.0\,GB)
    for a total of \textbf{10.3\,GB} --- 2.3$\times$ the MLX-only footprint.
    This is an important practical constraint for mobile deployment.
\end{itemize}

\subsection{Scaling Law: Larger Models Amortize Dispatch Overhead}
\label{sec:scaling}

Our 0.8B~vs.~2B comparison reveals a clear scaling trend:

\textbf{Model size $\uparrow$ $\Rightarrow$ dispatch overhead $\downarrow$}:
0.8B (hidden\_size=1024) incurs $\approx$250\,ms CoreML dispatch overhead at seq64;
2B (hidden\_size=1536) incurs \emph{zero} overhead --- TTFT exactly matches GPU baseline.
The cause is straightforward: larger hidden dimensions increase the
compute-per-dispatch ratio, making the fixed IPC/dispatch cost negligible.

\textbf{Prompt length $\uparrow$ $\Rightarrow$ prefill throughput $\uparrow$}:
This holds for both models.
0.8B achieves 4{,}128\,tok/s at seq512 (410 tokens) but only 22\,tok/s at
seq64 (6 tokens), because the ANE processes a fixed-size tensor and
longer real-token sequences waste less padding.

\textbf{Practical constraints of scaling up}:
\begin{itemize}[topsep=2pt,itemsep=1pt]
  \item \textbf{Conversion time}: grows super-linearly.
    2B seq512 requires $\approx$120\,min (vs.\ $\approx$22\,min for 0.8B
    seq512) on M2 Ultra.
  \item \textbf{First-load ANE compilation}: 2B seq256 takes 878\,s
    ($\approx$15\,min) on first load; cached thereafter.
  \item \textbf{Storage}: hybrid deployment stores both MLX weights and
    CoreML packages (e.g.\ 4.0\,GB + 3$\times$2.3\,GB = 10.9\,GB for 2B).
\end{itemize}

\textbf{Conclusion}: hybrid CoreML+MLX inference is most advantageous for
\emph{larger models with longer prompts}.
On mobile Apple Silicon (A-series, 6--8 GPU cores), the crossover point
falls even further, making ANE-assisted prefill beneficial for nearly all
conversational prompts (Table~\ref{tab:crossover}).

\subsection{Extrapolation to Other Apple Silicon Generations}
\label{sec:extrapolation}

The crossover point depends on GPU prefill throughput, which scales
approximately with GPU core count and memory bandwidth.
Table~\ref{tab:crossover} extrapolates from our M2 Ultra measurement to
other Apple Silicon chips, assuming ANE throughput scales proportionally
with TOPS and CoreML dispatch overhead remains roughly constant
($\approx$100\,ms for seq512).
\emph{These are estimates, not measured values.}

\begin{table}[h]
\centering
\footnotesize
\resizebox{\columnwidth}{!}{%
\begin{tabular}{@{}lrrrr@{}}
\toprule
Chip & GPU & Prefill (tok/s) & Crossover & Note \\
\midrule
M2 Ultra & 76c & 4{,}128 & $\sim$410 tok & measured \\
M1 Max   & 32c & $\sim$2{,}100 & $\sim$200 tok & est. \\
M1       &  8c & $\sim$500    & $\sim$50 tok  & est. \\
A17 Pro  &  6c & $\sim$400    & $\sim$40 tok  & est. \\
\bottomrule
\end{tabular}}
\caption{Estimated CoreML hybrid crossover point (prompt length above
  which ANE prefill outperforms GPU prefill) across Apple Silicon.
  GPU throughput linearly extrapolated from M2 Ultra; actual values
  vary with memory bandwidth and microarchitecture.}
\label{tab:crossover}
\end{table}

On mobile Apple Silicon (A17 Pro and similar), where GPU core count is
$\approx$10$\times$ lower than M2 Ultra, the estimated crossover falls
below 64 tokens --- meaning nearly every conversational prompt would
benefit from ANE-assisted prefill.
This suggests that Hybrid CoreML+MLX inference is most impactful on
mobile devices, where GPU resources are most constrained.

\subsection{Power Efficiency: ANE vs.\ GPU Prefill}
\label{sec:power}

\textbf{Theoretical argument.}
Prefill is a compute-intensive, fixed-shape workload: ideal for the ANE,
which is a purpose-built matrix-multiply accelerator operating at
$\approx$0.5--2\,W~\cite{anelm}.
The GPU performing the same operation draws $\approx$5--15\,W,
as it must maintain shader core clocks and memory controller throughput.
On mobile devices, sustained GPU compute during prefill triggers thermal
throttling, reducing subsequent decode speed.
The ANE does not share the GPU thermal budget, so offloading prefill to
the ANE preserves full GPU bandwidth for the decode phase.

\textbf{Empirical support from our data.}
Our ANE-LM Hybrid results confirm that moving prefill to the ANE does
\emph{not} degrade GPU decode performance:
measured decode throughput (67--70\,tok/s) is within noise of the MLX
GPU baseline (69--72\,tok/s, see Table~\ref{tab:pipeline_ttft}).
Since GPU behaviour during decode is identical whether prefill was handled
by the ANE or the GPU, the GPU thermal budget is entirely conserved for
the bandwidth-bound decode phase.

\textbf{Power measurements (M2 Ultra, Qwen3.5-0.8B FP16).}
We measure per-phase GPU, ANE, and CPU power using \texttt{powermetrics}
at 100\,ms sampling intervals.
Because MLX GPU prefill completes in $\approx$96\,ms per pass, we loop
prefill continuously for 15\,s to obtain a sustained measurement.
Table~\ref{tab:power} reports averaged readings.

\begin{table}[h]
\centering
\footnotesize
\resizebox{\columnwidth}{!}{%
\begin{tabular}{@{}lrrr@{}}
\toprule
Phase & GPU & ANE & CPU \\
\midrule
MLX GPU prefill (loop$^\dagger$) & 62.05\,W & 0.00\,W & 2.34\,W \\
MLX GPU decode                   & 14.17\,W & 0.00\,W & 2.98\,W \\
ANE-LM prefill                   &  0.22\,W & 1.58\,W & 3.77\,W \\
ANE-LM Hybrid decode             & 12.37\,W & 0.00\,W & 5.85\,W \\
ANE-LM pure                      &  0.16\,W & 1.42\,W & 5.47\,W \\
\bottomrule
\end{tabular}}
\caption{Average power per inference phase, M2 Ultra, Qwen3.5-0.8B FP16.
  \texttt{powermetrics} at 100\,ms intervals.
  $^\dagger$MLX prefill looped for 15\,s (332 iters) for steady-state.}
\label{tab:power}
\end{table}

The results quantify the thermal advantage of ANE-offloaded prefill:
GPU power drops from 62.05\,W to 0.22\,W --- a \textbf{282$\times$ reduction} ---
while ANE draws only 1.58\,W for the same computation.
The net reduction in package power is \textbf{60.25\,W}.
Decode GPU power is 14.17\,W (MLX baseline) vs.\ 12.37\,W (ANE-LM Hybrid),
a difference of 1.80\,W that falls within run-to-run variance,
confirming that the cache bridge introduces no thermal overhead in the
decode phase.

On mobile Apple Silicon (A-series, nominal TDP $\approx$3--8\,W total),
a 62\,W GPU prefill burst is physically impossible; thermal throttling
would activate within milliseconds.
In contrast, ANE prefill at $\approx$2\,W total chip draw remains well
within the mobile thermal envelope, enabling sustained inference without
clock-speed degradation.

\subsection{Hardware Evolution: M1--M4 vs.\ M5}

With M5, Apple integrated Neural Accelerators directly into GPU shader
cores, accessible via the Metal\,4 Tensor API without the CoreML dispatch
layer~\cite{apple_m5_mlx}.
This represents the hardware convergence our software pipeline approximates:

\begin{table}[h]
\centering
\footnotesize
\begin{tabular}{@{}p{1.2cm}p{3.2cm}p{1.5cm}@{}}
\toprule
Gen. & Prefill path & ANE access \\
\midrule
M1--M4 & CoreML $\to$ ANE (IPC daemon) & Indirect \\
M5+ & MLX $\to$ Metal Neural Acc. & Direct \\
\bottomrule
\end{tabular}
\caption{Evolution of ANE access across Apple Silicon generations.}
\label{tab:m5}
\end{table}

Our hybrid approach is most relevant for M1--M4 hardware with
long-context workloads where CoreML dispatch overhead is amortized.

% ══════════════════════════════════════════════════════════════════════════════
\section{Related Work}
\label{sec:related}

\textbf{MLX}~\cite{mlx}: GPU-focused array framework with lazy evaluation
and automatic differentiation, designed for Apple Silicon.

\textbf{CoreML / coremltools}: Apple's model deployment framework, routing
fixed-shape operations to ANE.
Used for prefill in our Hybrid pipeline.

\textbf{Anemll}~\cite{anemll}: Per-layer CoreML conversion for LLMs,
splitting each linear layer into a separate CoreML model.
Supports Llama2-family models.

\textbf{ANE}~\cite{ane}: Research repository demonstrating training on ANE
via private APIs for Stories110M-scale models.

\textbf{ANE-LM}~\cite{anelm}: Inference engine using private
\texttt{AppleNeuralEngine.framework} APIs for direct ANE dispatch,
supporting Qwen3 and Qwen3.5.

\textbf{llama.cpp}~\cite{llamacpp}: Portable LLM inference with CPU/Metal
backends, widely benchmarked on Apple Silicon.

% ══════════════════════════════════════════════════════════════════════════════
\section{Future Work}
\label{sec:future}

\begin{itemize}[topsep=2pt,itemsep=1pt]
  \item \textbf{2B-BF16 long context and 9B hybrid benchmark}: 2B-BF16 short
    and medium prompts now benchmarked (Table~\ref{tab:hybrid_2b}); seq512
    and 9B conversions are in progress.
  \item \textbf{seq\_len $>$ 512}: Longer context would further reduce the
    relative CoreML dispatch overhead.
  \item \textbf{Mobile validation}: Crossover-point estimates in
    Table~\ref{tab:crossover} are extrapolated; direct measurement on
    A-series hardware (iPhone, iPad) would validate the mobile claims.
  \item \textbf{Streaming cache bridge}: The current ANE-LM cache bridge
    serializes the full cache to disk then deserializes it; a zero-copy
    shared-memory approach would reduce bridge latency.
  \item \textbf{M5 validation}: Direct comparison between our hybrid
    approach and M5's Metal Neural Accelerator path would quantify
    the software--hardware equivalence.
\end{itemize}

% ══════════════════════════════════════════════════════════════════════════════
\section*{Acknowledgements}

This work builds on
ANE~\cite{ane} (private API research),
ANE-LM~\cite{anelm} (private API inference engine),
Anemll~\cite{anemll} (per-layer CoreML conversion),
MLX~\cite{mlx} (Apple's array framework), and
mlx-vlm~\cite{mlxvlm} (VLM support for MLX).

% ══════════════════════════════════════════════════════════════════════════════
\begin{thebibliography}{9}

\bibitem{apple_m5_mlx}
Apple ML Research.
\textit{Exploring LLMs on MLX with M5.}
Apple Machine Learning Research Blog, 2025.\\
\url{https://machinelearning.apple.com/research/exploring-llms-mlx-m5}

\bibitem{mlx}
Apple Inc.
\textit{MLX: An array framework for Apple Silicon.}
GitHub, 2023.
\url{https://github.com/ml-explore/mlx}

\bibitem{ane}
maderix.
\textit{ANE: Training neural networks on Apple Neural Engine.}
GitHub, 2024.
\url{https://github.com/maderix/ANE}

\bibitem{anemll}
Anemll.
\textit{Anemll: Apple Neural Engine ML framework.}
GitHub, 2024.
\url{https://github.com/Anemll/Anemll}

\bibitem{anelm}
AtomGradient.
\textit{ANE-LM: LLM inference on Apple Neural Engine via private API.}
GitHub, 2025.
\url{https://github.com/AtomGradient/ANE-LM}

\bibitem{llamacpp}
ggml-org.
\textit{llama.cpp: LLM inference in C/C++.}
GitHub, 2023.
\url{https://github.com/ggerganov/llama.cpp}

\bibitem{mlxvlm}
Blaizzy.
\textit{mlx-vlm: Vision language models for Apple Silicon.}
GitHub, 2024.
\url{https://github.com/Blaizzy/mlx-vlm}

\bibitem{qwen35}
Qwen Team.
\textit{Qwen3.5 Technical Report.}
Alibaba DAMO Academy, 2025.

\end{thebibliography}

\end{document}
